\textbf{Выберите один вариант лабораторной работы №4.} Здесь вы вручную
создаёте программу вычисления факториала на C, собираете её, запускаете
и отлаживаете. Вариант 2 --- усложнённая работа с загрузочным сектором и
журналом оценок; выполнять оба варианта не требуется.

Среда: Debian актуальной версии или Ubuntu с GCC, GDB, nano и Vim.
Используйте обычную учётную запись. Если нужна помощь с входом в Ubuntu,
терминалом и каталогами, откройте дополнительную методичку «Загрузка
Ubuntu и первые шаги в терминале».

\hypertarget{ux43fux43eux440ux44fux434ux43eux43a-ux441ux434ux430ux447ux438-ux43bux430ux431ux43eux440ux430ux442ux43eux440ux43dux44bux445-ux440ux430ux431ux43eux442--ux432ux430ux440ux438ux430ux43dux442-1}{%
\subsubsection{Порядок сдачи лабораторных работ --- вариант
1}\label{ux43fux43eux440ux44fux434ux43eux43a-ux441ux434ux430ux447ux438-ux43bux430ux431ux43eux440ux430ux442ux43eux440ux43dux44bux445-ux440ux430ux431ux43eux442--ux432ux430ux440ux438ux430ux43dux442-1}}

\begin{enumerate}
\def\labelenumi{\arabic{enumi}.}
\tightlist
\item
  \textbf{На занятии} выполните часть работы, запишите команды, расчёты
  и наблюдаемый результат и покажите его преподавателю. Если в процессе
  работы возникли вопросы, задайте их преподавателю.
\item
  \textbf{К следующей лабораторной} завершите работу и подготовьте
  отчёт: что сделали, как проверили, что ожидали и что наблюдали.
  Изучите использованные команды и понятия; подготовьтесь объяснить,
  как работает ваша программа, применённые механизмы и устройства компьютера.
  Сохраните исходники и результаты, чтобы повторить демонстрацию.
  \textbf{Отправьте преподавателю отчёт в формате PDF или HTML.}
\item
  \textbf{Во время следующей лабораторной} индивидуально защитите
  предыдущую работу: покажите преподавателю отчёт и работающий
  результат, объясните свои действия и ответьте на вопросы. Преподаватель
  сообщит, если нужно что-то исправить.
\item
  \textbf{На последнем занятии} защитите предыдущую лабораторную работу
  и сдайте оставшиеся задолженности. Преподаватель объявит текущую
  накопленную оценку за лабораторные работы.
\end{enumerate}

На семинаре №4 покажите текущую версию \texttt{factorial.c}, сборку или
отладку и запишите достигнутый результат. К семинару №5 завершите
выбранный вариант, сохраните исходник и отчёт для защиты.

\hypertarget{ux43fux43eux434ux433ux43eux442ux43eux432ux44cux442ux435-ux438ux43dux441ux442ux440ux443ux43cux435ux43dux442ux44b-ux438-ux43aux430ux442ux430ux43bux43eux433-ux43fux440ux43eux435ux43aux442ux430}{%
\subsubsection{Подготовьте инструменты и каталог
проекта}\label{ux43fux43eux434ux433ux43eux442ux43eux432ux44cux442ux435-ux438ux43dux441ux442ux440ux443ux43cux435ux43dux442ux44b-ux438-ux43aux430ux442ux430ux43bux43eux433-ux43fux440ux43eux435ux43aux442ux430}}

В терминале проверьте инструменты: \texttt{gcc\ -\/-version},
\texttt{gdb\ -\/-version}, \texttt{nano\ -\/-version},
\texttt{vim\ -\/-version}. Если чего-то нет и у вас есть право
устанавливать пакеты, в Debian/Ubuntu выполните
\texttt{sudo\ apt\ update} и
\texttt{sudo\ apt\ install\ gcc\ libc6-dev\ gdb\ nano\ vim}. На другом
дистрибутиве названия пакетов могут отличаться. Команды ниже вводятся
без знака \texttt{\$} приглашения оболочки.

\begin{verbatim}
mkdir -p ~/aiya/lab04-factorial
cd ~/aiya/lab04-factorial
pwd
ls
\end{verbatim}

\texttt{mkdir\ -p} создаёт вложенные каталоги и не выдаёт ошибку, если
они уже существуют; \texttt{\textasciitilde{}} --- ваш домашний каталог,
\texttt{cd} меняет текущий каталог. \texttt{pwd} должен оканчиваться на
\texttt{/aiya/lab04-factorial}. \texttt{ls} показывает файлы в нём: в
первый раз каталог пуст.

\hypertarget{ux43dux430ux43fux438ux448ux438ux442ux435-ux43fux435ux440ux432ux443ux44e-ux432ux435ux440ux441ux438ux44e-ux432-nano}{%
\subsubsection{Напишите первую версию в
nano}\label{ux43dux430ux43fux438ux448ux438ux442ux435-ux43fux435ux440ux432ux443ux44e-ux432ux435ux440ux441ux438ux44e-ux432-nano}}

Введите \texttt{nano\ factorial.c} и наберите код \textbf{вручную}, не
копируя готовый файл:

\begin{verbatim}
#include <stdio.h>

unsigned long long factorial(unsigned int n)
{
    unsigned long long result = 1;
    for (unsigned int i = 1; i <= n; ++i) {
        result *= i;
    }
    return result;
}

int main(void)
{
    int n = 5;
    unsigned long long value = factorial((unsigned int)n);
    printf("factorial(%d) = %llu\n", n, value);
    return 0;
}
\end{verbatim}

Сохраните файл клавишами \textbf{Ctrl+O}, \textbf{Enter}, выйдите через
\textbf{Ctrl+X}. Внизу nano запись \texttt{\^{}O} означает Ctrl+O.
Проверьте наличие файла через \texttt{ls} и при желании прочитайте его
через \texttt{cat\ factorial.c}.

\texttt{main} --- начало программы. \texttt{factorial} умножает числа от
1 до \texttt{n}; переменная \texttt{result} начинается с 1, поэтому при
\texttt{n\ =\ 0} цикл не выполнится и получится \texttt{0!\ =\ 1}.
\texttt{unsigned\ long\ long} хранит беззнаковое целое большой
разрядности. \texttt{printf} печатает в стандартный вывод: \texttt{\%d}
соответствует \texttt{int}, \texttt{\%llu} ---
\texttt{unsigned\ long\ long}, \texttt{\textbackslash{}n} завершает
строку. После исполнения \texttt{return\ 0} сообщает оболочке об
успешном завершении.

\hypertarget{ux441ux43eux431ux435ux440ux438ux442ux435-ux438-ux437ux430ux43fux443ux441ux442ux438ux442ux435-ux43fux440ux43eux433ux440ux430ux43cux43cux443}{%
\subsubsection{Соберите и запустите
программу}\label{ux441ux43eux431ux435ux440ux438ux442ux435-ux438-ux437ux430ux43fux443ux441ux442ux438ux442ux435-ux43fux440ux43eux433ux440ux430ux43cux43cux443}}

\begin{verbatim}
gcc -std=c11 -Wall -Wextra -Wpedantic -g -O0 \
  factorial.c -o factorial
./factorial
echo $?
\end{verbatim}

Ожидаемый вывод: \texttt{factorial(5)\ =\ 120}, затем \texttt{0} от
\texttt{echo\ \$?} (смотрите код завершения \textbf{сразу после} запуска
программы). Если \texttt{gcc} напечатал ошибку, найдите строку и
исправьте код; при неуспешной сборке \textbf{не запускайте старый} файл
\texttt{factorial}. Оболочка ищет программы в каталогах \texttt{PATH},
поэтому исполняемый файл в текущем каталоге запускаем через
\texttt{./factorial}.

\begin{longtable}[]{@{}
  >{\raggedright\arraybackslash}p{(\columnwidth - 2\tabcolsep) * \real{0.2700}}
  >{\raggedright\arraybackslash}p{(\columnwidth - 2\tabcolsep) * \real{0.6900}}@{}}
\toprule
\begin{minipage}[b]{\linewidth}\raggedright
Флаг GCC
\end{minipage} & \begin{minipage}[b]{\linewidth}\raggedright
Что делает
\end{minipage} \\
\midrule
\endhead
\texttt{-std=c11} & Выбирает стандарт языка C11. \\
\texttt{-Wall} & Включает распространённые предупреждения. \\
\texttt{-Wextra} & Включает дополнительные предупреждения. \\
\texttt{-Wpedantic} & Предупреждает об отклонениях от выбранного
стандарта языка. \\
\texttt{-g} & Добавляет информацию для отладчика: имена функций и строки
исходника. \\
\texttt{-O0} & Выключает оптимизацию, чтобы проще сопоставлять код с
шагами GDB. \\
\texttt{-o\ factorial} & Задаёт имя исполняемого файла; без него GCC
обычно создаёт \texttt{a.out}. \\
\bottomrule
\end{longtable}

Позже можно попробовать \texttt{-Werror} (считать предупреждения
ошибками), \texttt{-O2} (оптимизировать программу; при отладке
отображение переменных может измениться) и \texttt{-c} (получить
объектный файл без линковки и готовой программы). Для основного опыта
оставьте \texttt{-g\ -O0}.

Файл \texttt{factorial.c} --- текст исходника на диске;
\texttt{factorial} после успешной сборки --- исполняемый файл на диске.
При \texttt{./factorial} Linux создаёт \textbf{процесс}. После изменения
\texttt{.c} старая исполняемая программа не меняется, пока вы не
повторите сборку GCC.

\hypertarget{ux438ux437ux43cux435ux43dux438ux442ux435-ux43fux440ux43eux433ux440ux430ux43cux43cux443-ux447ux435ux440ux435ux437-vim-ux438-ux434ux43eux431ux430ux432ux44cux442ux435-ux432ux432ux43eux434}{%
\subsubsection{Измените программу через Vim и добавьте
ввод}\label{ux438ux437ux43cux435ux43dux438ux442ux435-ux43fux440ux43eux433ux440ux430ux43cux43cux443-ux447ux435ux440ux435ux437-vim-ux438-ux434ux43eux431ux430ux432ux44cux442ux435-ux432ux432ux43eux434}}

Сначала посмотрите редактор на знакомом файле:
\texttt{vim\ factorial.c}. Vim открывается в \textbf{обычном режиме}:
нажмите \texttt{i} для режима вставки и исправьте \texttt{int\ n\ =\ 5;}
на \texttt{int\ n\ =\ 6;}. Нажмите \textbf{Esc}, затем наберите
\texttt{:wq} и \textbf{Enter} --- сохранить и выйти. Пересоберите и
запустите: ожидается \texttt{factorial(6)\ =\ 720}. Полезные команды из
обычного режима: \texttt{:w} сохраняет без выхода, \texttt{:q} выходит
при отсутствии несохранённых правок, \texttt{:q!} выходит \textbf{без
сохранения}. Если буквы не вводятся в файл, вы остались в обычном
режиме: нажмите \texttt{i}.

Теперь снова откройте \texttt{vim\ factorial.c} (или
\texttt{nano\ factorial.c}) и \textbf{замените содержимое} следующим
кодом. Он читает одно целое число и явно разделяет обычный результат и
ошибки:

\begin{verbatim}
#include <stdio.h>

unsigned long long factorial(unsigned int n)
{
    unsigned long long result = 1;
    for (unsigned int i = 1; i <= n; ++i) {
        result *= i;
    }
    return result;
}

int main(void)
{
    int n;
    if (fscanf(stdin, "%d", &n) != 1) {
        fprintf(stderr, "Ошибка ввода.\n");
        return 1;
    }
    if (n < 0 || n > 20) {
        fprintf(stderr, "Нужно 0..20.\n");
        return 1;
    }
    unsigned long long value = factorial((unsigned int)n);
    printf("factorial(%d) = %llu\n", n, value);
    return 0;
}
\end{verbatim}

Сохраните изменения и \textbf{снова выполните команду GCC}. Запустите
\texttt{./factorial}, введите \texttt{5} и нажмите Enter: результат
должен быть \texttt{factorial(5)\ =\ 120}. Проверьте \texttt{0},
\texttt{1}, \texttt{10}, \texttt{20}, \texttt{-1}, \texttt{21}, буквы
вместо числа и пустой ввод (для конца ввода в интерактивном терминале
нажмите \textbf{Ctrl+D} на пустой строке). Сравните с таблицей:

\begin{longtable}[]{@{}
  >{\raggedright\arraybackslash}p{(\columnwidth - 2\tabcolsep) * \real{0.2700}}
  >{\raggedright\arraybackslash}p{(\columnwidth - 2\tabcolsep) * \real{0.6900}}@{}}
\toprule
\begin{minipage}[b]{\linewidth}\raggedright
Ввод
\end{minipage} & \begin{minipage}[b]{\linewidth}\raggedright
Ожидаемое поведение
\end{minipage} \\
\midrule
\endhead
\texttt{0}, \texttt{1} & \texttt{factorial(0)\ =\ 1},
\texttt{factorial(1)\ =\ 1} \\
\texttt{5}, \texttt{10} & \texttt{120}, \texttt{3628800}
соответственно \\
\texttt{20} & \texttt{2432902008176640000} \\
\texttt{-1}, \texttt{21}, буквы, конец ввода & Сообщение об ошибке,
ненулевой код завершения \\
\bottomrule
\end{longtable}

Число \texttt{20!} помещается в обычный 64-битный
\texttt{unsigned\ long\ long}, а \texttt{21!\ =\ 51090942171709440000}
больше максимально представимого значения \texttt{18446744073709551615};
поэтому программа ограничивает ввод. Данный учебный пример читает
\textbf{первое целое число типа \texttt{int}}, а не проверяет всю
строку: после \texttt{5} лишний текст не распознаётся как ошибка. Ввод
числа вне диапазона \texttt{int} тоже не рассчитан на обработку этим
простым вариантом \texttt{fscanf}; не выдавайте такую проверку за полную
защиту от любого ввода.

\hypertarget{ux444ux430ux439ux43bux44b-ux438-ux442ux440ux438-ux441ux442ux430ux43dux434ux430ux440ux442ux43dux44bux445-ux43fux43eux442ux43eux43aux430}{%
\subsubsection{Файлы и три стандартных
потока}\label{ux444ux430ux439ux43bux44b-ux438-ux442ux440ux438-ux441ux442ux430ux43dux434ux430ux440ux442ux43dux44bux445-ux43fux43eux442ux43eux43aux430}}

Файл --- последовательность байтов, сохраняемая файловой системой.
Текстовый исходник, исполняемый файл и \texttt{input.txt} --- разные
файлы; расширение \texttt{.c} само по себе не делает файл исполняемым.
Программа на C работает с потоком \texttt{FILE\ *} (абстракция
стандартной библиотеки), который может читать из терминала или файла и
писать в терминал или файл. В Linux у процесса при запуске обычно
открыты стандартные дескрипторы: 0 (ввод), 1 (вывод), 2 (ошибки).
Объекты C \texttt{stdin}, \texttt{stdout}, \texttt{stderr} связаны с
этими потоками, но \texttt{FILE\ *} --- \textbf{не номер дескриптора}.

\texttt{fscanf(stdin,\ "\%d",\ \&n)} читает число из стандартного ввода
и возвращает количество успешно прочитанных значений (здесь нужно
\texttt{1}). Упрощённая запись --- \texttt{scanf("\%d",\ \&n)}; в обоих
случаях нужен адрес \texttt{\&n} и проверка результата.
\texttt{printf(...)} пишет в \texttt{stdout}, то же можно записать
\texttt{fprintf(stdout,\ ...)}. \texttt{fprintf(stderr,\ ...)} пишет в
отдельный поток ошибок. Эти функции \textbf{выводят} данные; для ввода
служат \texttt{scanf}/\texttt{fscanf}.

Скопируйте \texttt{input.txt} из комплекта материалов в рабочий каталог
или создайте его командой \texttt{nano\ input.txt} и запишите число
\texttt{5}. Если вы работаете без архива, достаточно варианта с nano.
Затем выполните:

\begin{verbatim}
./factorial < input.txt > output.txt 2> errors.txt
echo $?
cat output.txt
cat errors.txt
\end{verbatim}

\texttt{\textless{}} подаёт файл в \texttt{stdin},
\texttt{\textgreater{}} сохраняет \texttt{stdout} (при повторе
\textbf{перезаписывает} файл), \texttt{2\textgreater{}} сохраняет
\texttt{stderr}. Ожидайте ответ в \texttt{output.txt} и пустой
\texttt{errors.txt}. Измените \texttt{input.txt} на \texttt{21},
повторите запуск и \textbf{сразу} проверьте \texttt{echo\ \$?}: теперь
\texttt{output.txt} пуст, а сообщение --- в \texttt{errors.txt}.
\texttt{\textgreater{}\textgreater{}} добавляет вывод к существующему
файлу вместо перезаписи. Перенаправления оболочки влияют на запущенную
программу; диагностические сообщения самого GCC при сборке --- отдельный
запуск другого процесса.

\hypertarget{ux43dux430ux439ux434ux438ux442ux435-ux43eux448ux438ux431ux43aux443-ux432-gdb}{%
\subsubsection{Найдите ошибку в
GDB}\label{ux43dux430ux439ux434ux438ux442ux435-ux43eux448ux438ux431ux43aux443-ux432-gdb}}

Восстановите в \texttt{input.txt} число \texttt{5} и \textbf{снова
соберите} программу с \texttt{-g\ -O0}: \texttt{-g} сохраняет имена
переменных и соответствие строк исходнику, а \texttt{-O0} не даёт
оптимизатору переставить или убрать интересующие нас вычисления.
Работайте из каталога с \texttt{factorial.c}, \texttt{factorial} и
\texttt{input.txt}. Сначала в \textbf{обычной оболочке} наберите:

\begin{verbatim}
gdb -q ./factorial
\end{verbatim}

\texttt{-q} скрывает вводный баннер. Появится приглашение
\textbf{\texttt{(gdb)}}; дальнейшие команды вводите уже \emph{в
отладчике}, без текста \texttt{(gdb)}. Выражение \texttt{./factorial}
--- путь к исполняемому файлу, а не команда внутри GDB. На остановке GDB
показывает строку, которую программа \textbf{собирается выполнить}: до
первого \texttt{next} значение ещё не инициализированной локальной
переменной читать нельзя. Номера строк ниже относятся к показанному выше
коду; если вы добавляли строки в \texttt{factorial.c}, номера сдвинутся.

\hypertarget{ux448ux430ux433-1-ux43eux441ux442ux430ux43dux43eux432ux438ux442ux435ux441ux44c-ux432-ux444ux443ux43dux43aux446ux438ux438}{%
\paragraph{Шаг 1. Остановитесь в
функции}\label{ux448ux430ux433-1-ux43eux441ux442ux430ux43dux43eux432ux438ux442ux435ux441ux44c-ux432-ux444ux443ux43dux43aux446ux438ux438}}

\begin{verbatim}
break factorial
info breakpoints
run < input.txt
print n
info args
backtrace
\end{verbatim}

\texttt{break\ factorial} ставит точку останова на входе в функцию. GDB
сообщает её \textbf{номер} (\texttt{Breakpoint\ 1}) и строку.
\texttt{info\ breakpoints} показывает номера, состояние и места всех
остановок; используйте реальные номера, которые видите у себя. Можно
ставить остановку и по строке, например \texttt{break\ factorial.c:7}
--- это строка \texttt{result\ *=\ i;} в исходнике из методички. Не
добавляйте её пока, чтобы не изменить следующий опыт.

\texttt{run\ \textless{}\ input.txt} \textbf{начинает новый запуск}
программы: \texttt{\textless{}} подаёт файл в её стандартный ввод.
Отладчик остановится в \texttt{factorial} до выполнения строки
\texttt{unsigned\ long\ long\ result\ =\ 1;}. В сообщении
\texttt{Breakpoint\ ...\ factorial\ (n=5)\ ...} важны функция, аргумент,
имя файла и текущая строка; адреса и номера внутренних кадров могут
отличаться. Если забыть \texttt{\textless{}\ input.txt}, программа может
ждать число в терминале: введите \texttt{5} и Enter либо прервите
выполнение клавишами \textbf{Ctrl+C} и повторите
\texttt{run\ \textless{}\ input.txt}.

\texttt{print\ n} покажет значение аргумента \texttt{5}. Запись вида
\texttt{\$1\ =\ 5} --- это номер \textbf{результата команды
\texttt{print} в истории GDB}, а не переменная C и не код завершения
\texttt{\$?} из оболочки. \texttt{info\ args} перечислит аргументы
текущей функции. \texttt{backtrace} (или \texttt{bt}) покажет цепочку
вызовов: \texttt{\#0\ factorial} --- где мы остановились,
\texttt{\#1\ main} --- кто вызвал функцию. Команда \texttt{frame\ 1}
выбирает кадр \texttt{main} для просмотра его переменных,
\texttt{frame\ 0} возвращает контекст просмотра в \texttt{factorial};
сами эти команды \textbf{не выполняют программу}. Находясь в
\texttt{main}, можно посмотреть его \texttt{n}, а \texttt{result} там
недоступна.

\hypertarget{ux448ux430ux433-2-ux432ux44bux43fux43eux43bux43dux44fux439ux442ux435-ux441ux442ux440ux43eux43aux438}{%
\paragraph{Шаг 2. Выполняйте
строки}\label{ux448ux430ux433-2-ux432ux44bux43fux43eux43bux43dux44fux439ux442ux435-ux441ux442ux440ux43eux43aux438}}

Убедитесь, что выбран \texttt{frame\ 0}, затем введите:

\begin{verbatim}
next
print result
info locals
next
print i
next
print result
\end{verbatim}

Первый \texttt{next} выполнит \texttt{result\ =\ 1} и поставит указатель
на заголовок \texttt{for}; теперь \texttt{print\ result} выдаёт
\texttt{1}. \texttt{info\ locals} показывает доступные локальные
переменные, но на заголовке \texttt{for} значение \texttt{i} до
инициализации не следует считать результатом вычисления. Следующий
\texttt{next} войдёт в первую итерацию и остановится перед
\texttt{result\ *=\ i}: здесь \texttt{print\ i} покажет \texttt{1}. Ещё
один \texttt{next} выполнит умножение \texttt{1\ ×\ 1}, поэтому
\texttt{result} всё ещё \texttt{1}. У цикла есть заголовок, тело и
переход к следующей итерации: один \texttt{next} \textbf{не обязательно
равен одной итерации}. Если вы оказались на другой строке, посмотрите на
строку, которую показывает GDB, и выполните следующую команду осознанно.

\texttt{next} выполняет текущую строку, \textbf{не заходя внутрь}
вызываемой функции; \texttt{step} на строке с вызовом заходит в неё. Для
отдельного сравнения вернитесь в \texttt{main} новым запуском:
\texttt{break\ main}, \texttt{run\ \textless{}\ input.txt}, затем
делайте \texttt{next} до строки \texttt{value\ =\ factorial(...)} и
введите \texttt{step} --- вы попадёте в \texttt{factorial}.
\texttt{finish} продолжит выполнение текущей функции до \texttt{return}
и остановится в вызывающей функции, показав возвращённое значение. При
\texttt{step} на строке \texttt{printf} можно попасть в код библиотеки;
для первого опыта удобнее \texttt{next} или \texttt{finish}. Перед
следующим опытом удалите временную точку \texttt{break\ main} командой
\texttt{delete\ НОМЕР} (номер смотрите через
\texttt{info\ breakpoints}). При каждом повторном \texttt{run}
выполнение начинается заново, но остальные точки останова остаются до их
удаления.

\hypertarget{ux448ux430ux433-3-ux441ux43bux435ux434ux438ux442ux435-ux437ux430-ux43fux440ux43eux438ux437ux432ux435ux434ux435ux43dux438ux435ux43c}{%
\paragraph{Шаг 3. Следите за
произведением}\label{ux448ux430ux433-3-ux441ux43bux435ux434ux438ux442ux435-ux437ux430-ux43fux440ux43eux438ux437ux432ux435ux434ux435ux43dux438ux435ux43c}}

Проверьте через \texttt{info\ breakpoints}, что оставлена точка
\texttt{break\ factorial}, но нет временной остановки в \texttt{main}.
Начните новый запуск командой \texttt{run\ \textless{}\ input.txt}: при
остановке на входе в \texttt{factorial} сначала выполните
\textbf{\texttt{next}}, чтобы инициализировать \texttt{result}. Затем:

\begin{verbatim}
watch result
info breakpoints
continue
print i
print result
continue
continue
continue
print result
\end{verbatim}

\texttt{watch\ result} --- \textbf{точка наблюдения}: GDB остановится,
когда значение \texttt{result} \emph{изменится}, и покажет
\texttt{Old\ value} и \texttt{New\ value}. Первое умножение
\texttt{1\ ×\ 1} не меняет \texttt{result}, поэтому первой новой
величиной будет \textbf{2 при \texttt{i\ =\ 2}}. Затем повторные
\texttt{continue} приведут к изменениям \textbf{2 → 6 → 24 → 120}. На
последней остановке \texttt{print\ result} покажет \texttt{120};
следующее \texttt{continue} завершит программу с напечатанным ответом.
Не пытайтесь читать \texttt{result} уже после завершения: локальная
переменная существовала только во время вызова функции.

Для вывода значения на \textbf{каждой} остановке можно выполнить
\texttt{display\ result}: GDB присвоит этому отображению номер и будет
печатать значение автоматически, пока переменная доступна. Узнать номер
можно через \texttt{info\ display}, отключить ---
\texttt{undisplay\ НОМЕР}. \texttt{display} лишь показывает значение, а
\texttt{watch} останавливает выполнение при изменении. Список остановок
и их номера показывает \texttt{info\ breakpoints};
\texttt{disable\ НОМЕР} временно выключает точку, \texttt{enable\ НОМЕР}
включает обратно, \texttt{delete\ НОМЕР} удаляет её. Не подставляйте
номер наугад: в список входят и \texttt{break}, и \texttt{watch}.

\hypertarget{ux448ux430ux433-4-ux438ux441ux441ux43bux435ux434ux443ux439ux442ux435-ux441ux442ux435ux43a-ux438-ux43fux430ux43cux44fux442ux44c}{%
\paragraph{Шаг 4. Исследуйте стек и
память}\label{ux448ux430ux433-4-ux438ux441ux441ux43bux435ux434ux443ux439ux442ux435-ux441ux442ux435ux43a-ux438-ux43fux430ux43cux44fux442ux44c}}

После завершения предыдущего запуска локальная точка
\texttt{watch\ result} перестаёт действовать, поскольку переменная вышла
из области видимости. Снова запустите
\texttt{run\ \textless{}\ input.txt}, остановитесь в \texttt{factorial},
выполните \texttt{next} для инициализации \texttt{result} и поставьте
\texttt{watch\ result} заново. Четыре раза выполните \texttt{continue}:
остановитесь на изменении с \texttt{24} на \texttt{120}. Пока вы внутри
функции, выполните:

\begin{verbatim}
backtrace
info locals
print sizeof(result)
print &result
x/8xb &result
\end{verbatim}

\texttt{backtrace} снова покажет \texttt{\#0\ factorial} и
\texttt{\#1\ main}. \texttt{info\ locals} позволяет сопоставить
\texttt{i} и \texttt{result}. \texttt{sizeof(result)} в этой учебной
среде x86-64 равен \textbf{8 байтам}. \texttt{\&result} --- адрес
переменной в \emph{виртуальной} памяти процесса: конкретное число не
переписывайте в отчёт как универсальный адрес. В команде
\texttt{x/8xb\ \&result} буква \textbf{\texttt{x}} означает просмотр
памяти, \textbf{\texttt{8}} --- восемь элементов, следующий
\textbf{\texttt{x}} --- вывод в шестнадцатеричном виде,
\textbf{\texttt{b}} --- размер элемента в один байт. Для
\texttt{result\ =\ 120} (\texttt{0x78}) на типичном Linux x86-64
ожидаются байты \texttt{0x78\ 0x00\ 0x00\ 0x00\ 0x00\ 0x00\ 0x00\ 0x00}:
младший байт расположен первым (\emph{little-endian}). В начале опыта
при \texttt{result\ =\ 1} первым байтом будет \texttt{0x01}; не
смешивайте две разные остановки. Не путайте \texttt{x/8xb} со
стандартной командой оболочки или с выражением C.

\hypertarget{ux448ux430ux433-5-ux43dux430ux439ux434ux438ux442ux435-ux43eux448ux438ux431ux43aux443-ux432-ux443ux441ux43bux43eux432ux438ux438}{%
\paragraph{Шаг 5. Найдите ошибку в
условии}\label{ux448ux430ux433-5-ux43dux430ux439ux434ux438ux442ux435-ux43eux448ux438ux431ux43aux443-ux432-ux443ux441ux43bux43eux432ux438ux438}}

Выйдите из GDB командой \texttt{quit}. Через nano или Vim
\textbf{намеренно} измените условие в цикле на
\texttt{i\ \textless{}\ n} и пересоберите с \texttt{-g\ -O0}. Для ввода
\texttt{5} программа теперь выдаёт \textbf{24} вместо 120: последний
множитель не используется. Снова откройте GDB и остановитесь
непосредственно \textbf{перед умножением}:

\begin{verbatim}
break factorial.c:7
run < input.txt
print i
print result
continue
\end{verbatim}

Повторяйте \texttt{print\ i} и \texttt{continue} на остановках: до
строки \texttt{result\ *=\ i;} вы дойдёте при
\texttt{i\ =\ 1,\ 2,\ 3,\ 4}, но \textbf{не при \texttt{i\ =\ 5}}. Если
строки в вашем файле сдвинулись, поставьте точку по актуальному номеру
строки умножения (\texttt{list\ factorial} показывает исходник) или на
входе в \texttt{factorial} и идите через \texttt{next}. Запишите, какие
значения наблюдали и почему ответ равен 24. Восстановите
\texttt{i\ \textless{}=\ n}, \textbf{пересоберите} программу и повторите
проверку для \texttt{0}, \texttt{5} и \texttt{20}. Менять ожидаемый
ответ вместо исправления цикла нельзя.

Если GDB пишет \texttt{No\ symbol\ ...\ in\ current\ context},
проверьте, остановлена ли программа \textbf{внутри нужной функции}:
например, \texttt{i} видна только в области действия цикла, а после
завершения процесса локальные переменные недоступны. Сообщение
\texttt{The\ program\ is\ not\ being\ run} означает, что нужно выполнить
\texttt{run}. Если GDB не знает имён и строк исходника, проверьте сборку
с \texttt{-g} и то, что отлаживаете \textbf{заново собранный}
\texttt{./factorial}. Если вывод отладчика отличается, попробуйте
\texttt{help\ КОМАНДА} (например, \texttt{help\ watch}); \texttt{quit}
закрывает GDB, а \texttt{Ctrl+C} прерывает ожидающую ввода или
работающую программу.

\begin{longtable}[]{@{}
  >{\raggedright\arraybackslash}p{(\columnwidth - 2\tabcolsep) * \real{0.2700}}
  >{\raggedright\arraybackslash}p{(\columnwidth - 2\tabcolsep) * \real{0.6900}}@{}}
\toprule
\begin{minipage}[b]{\linewidth}\raggedright
Команда
\end{minipage} & \begin{minipage}[b]{\linewidth}\raggedright
Краткая памятка
\end{minipage} \\
\midrule
\endhead
\texttt{break}, \texttt{info\ breakpoints} & Остановиться в функции или
на строке, показать точки и их номера. \\
\texttt{run}, \texttt{continue} & Запустить сначала с вводом из файла;
продолжить текущий запуск. \\
\texttt{next}, \texttt{step}, \texttt{finish} & Выполнить строку без
захода в вызов, войти в функцию, завершить функцию. \\
\texttt{print}, \texttt{info\ args}, \texttt{info\ locals} & Узнать
значение выражения, аргументы и локальные переменные. \\
\texttt{watch}, \texttt{display} & Остановиться при изменении; показать
значение при остановках. \\
\texttt{backtrace}, \texttt{frame} & Показать стек вызовов и выбрать
кадр для просмотра. \\
\texttt{x/8xb}, \texttt{help}, \texttt{quit} & Прочитать восемь байтов
памяти, открыть справку, выйти из GDB. \\
\bottomrule
\end{longtable}

\hypertarget{ux441ux430ux43cux43eux441ux442ux43eux44fux442ux435ux43bux44cux43dux430ux44f-ux440ux430ux431ux43eux442ux430-ux438-ux43eux442ux447ux451ux442}{%
\subsubsection{Самостоятельная работа и
отчёт}\label{ux441ux430ux43cux43eux441ux442ux43eux44fux442ux435ux43bux44cux43dux430ux44f-ux440ux430ux431ux43eux442ux430-ux438-ux43eux442ux447ux451ux442}}

Добавьте внутри цикла после умножения строку
\texttt{fprintf(stderr,\ "i=\%u\ result=\%llu\textbackslash{}n",\ i,\ result);}
и пересоберите. Запустите программу с перенаправлениями: убедитесь, что
ответ остаётся в \texttt{output.txt}, а шаги вычисления попадают в
\texttt{errors.txt}. Объясните, почему для успеха код завершения
остаётся \texttt{0}, хотя в \texttt{stderr} появился текст. После опыта
удалите дополнительный диагностический вывод или оставьте его и укажите
это в отчёте.

В отчёте приведите команды создания каталога, сборки, запуска и GDB;
сравнение ожидаемого и наблюдаемого для \texttt{0}, \texttt{5},
\texttt{20}, \texttt{-1}, \texttt{21} и ошибочного ввода; состояние
\texttt{stdout}, \texttt{stderr} и кода завершения; найденную и
исправленную ошибку \texttt{i\ \textless{}\ n}; объяснение различий
между исходником, исполняемым файлом и процессом. Для ориентира по
структуре посмотрите «Пример отчёта: сумма двух чисел» в дополнительных
материалах; пример относится к другой, демонстрационной задаче.
Отправьте преподавателю \textbf{свой} отчёт в PDF или HTML. Сохраните
свой \texttt{factorial.c} и проверочные файлы отдельно для
индивидуальной защиты. Покажите преподавателю собственный результат и
объясните его, не только продемонстрируйте готовый вывод.
